\writeSectionFrame{\presentationTitle}{Fliegengewicht}
\begin{frame}{Steckbrief}
	\begin{itemize}
 		\item Deutscher Name: Fliegengewicht
 		\item Englischer Name: Flyweight
 		\item Gruppe: Strukturmuster
 	\end{itemize}
\end{frame}

\begin{frame}{Fliegengewicht}
	\begin{block}{Zweck}
 		Nutze Objekte kleinster Granularität gemeinsam, um große Mengen von ihnen effizient benutzen zu können.
 	\end{block}
\end{frame}

\begin{frame}{Allgemeines Klassendiagramm mit Vererbung}
    \begin{tikzpicture}[scale=0.7, transform shape]
            \umlclass[x=-5, y=0]{Flyweight}{
            +operation(extrinsicState)
            }{}
            
            \umlclass[x=0, y=3]{ConcreteFlyweight}{
            -intrinsicState
            }{
            +operation(extrinsicState)
            }
            
            \umlclass[x=0, y=-3]{UnsharedConcreteFlyweight}{
            -allState
            }{
            +operation(extrinsicState)
            }
        
            \umlclass[x=5, y=0]{FlyweightFactory}{
            -flyweights : Flyweight[]
            }{
            +getFlyweight(key) : Flyweight
            }
    
            \umlclass[x=5, y=3]{Client}{
            }{
            }
        
            \umlinherit[geometry=-|]{ConcreteFlyweight}{Flyweight}
            \umlinherit[geometry=-|]{UnsharedConcreteFlyweight}{Flyweight}
            \umlaggreg{FlyweightFactory}{Flyweight}
            \umluniassoc{Client}{FlyweightFactory}
    \end{tikzpicture}
\end{frame}
\note{
    Der intrinsische Zustand ist gemeinsam und wird von mehreren Objekten geteilt. Er ist unveränderlich und intern an das Fliegengewicht gebunden. Der extrinsische Zustand ist spezifisch für eine bestimmte Verwendung und wird außerhalb des Fliegengewicht-Objekts gespeichert. Er wird dem Fliegengewicht zur Laufzeit übergeben und ist veränderlich. Fliegengewicht-Objekte repräsentieren den intrinsischen Zustand und werden geteilt. Das Muster sorgt dafür, dass für gleiche intrinsische Zustände nur eine einzige Instanz existiert, wodurch Speicherplatz eingespart wird. Die FlyweightFactory verwaltet die Fliegengewicht-Objekte. Sie sorgt dafür, dass wiederverwendbare Fliegengewichte erstellt oder bereitgestellt werden, und stellt sicher, dass das richtige Fliegengewicht für eine bestimmte Anfrage verwendet wird. Der Client verwendet die FlyweightFactory, um Fliegengewicht-Objekte zu erhalten und diese mit extrinsischen Zuständen zu kombinieren, um die benötigten Operationen durchzuführen. UnsharedConcreteFlyweight speichert spezifische Zustände, die nicht geteilt werden.
}

\frameWithImageOnly{\BILDERPFAD/fliegengewicht1.jpg}
\note{
	Wir möchten gerne die Grundlage für ein kleines Spiel schaffen. Der Spieler bewegt sich über 
	eine Karte, die aus Kacheln besteht. Kacheln haben unterschiedliche Eigenschaften. Im Wald kann
	er Essen finden, kann aber auch angegriffen werden. In der Wüste verbraucht er Wasser und auf 
	dem Fluss kann er Wasser sammeln. Der Spieler verliert Gesundheit, wenn er angegriffen wird
	und wenn er in der Wüste kein Wasser mehr hat. Hat er keine Gesundheit mehr, ist er tot. 
	Dabei haben die einzelnen Kacheln, über die der Spieler laufen kann, keinen gemeinsamen 
	Zustand. 
}

\begin{frame}{Klassendiagramm - Karte für Spiel}
    \begin{tikzpicture}[scale=0.7, transform shape]
            \umlclass[x=-5, y=0]{TileFlyweight}{
            }{}
            
            \umlclass[x=0, y=3]{DesertTile}{
            }{
            }
    
            \umlclass[x=0, y=-1]{RiverTile}{
            }{
            }
    
            \umlclass[x=0, y=1]{ForestTile}{
            }{
            }
            
            \umlclass[x=0, y=-3]{UnsharedTile}{
            }{
            }
        
            \umlclass[x=5, y=0]{TileFlyweightFactory}{
            }{
            }
    
            \umlclass[x=5, y=3]{GameMap}{
            }{
            }
        
            \umlimpl[geometry=-|]{DesertTile}{TileFlyweight}
            \umlimpl{RiverTile}{TileFlyweight}
            \umlimpl{ForestTile}{TileFlyweight}
            \umlimpl[geometry=-|]{UnsharedTile}{TileFlyweight}
            \umlaggreg{TileFlyweightFactory}{TileFlyweight}
            \umluniassoc{GameMap}{TileFlyweightFactory}
    \end{tikzpicture}
\end{frame}

\begin{frame}[fragile]{Fliegengewicht}
	\begin{exampleblock}{TileFlyweight}
 		\begin{lstlisting}
public abstract class TileFlyweight {
	private final Texture texture;
	
	public TileFlyweight(Texture texture) {
		super();
		this.texture = texture;
	}

	public abstract void walk(Player player);
	public abstract String render();
	
	@Override
	public String toString() {
		return texture.toString() + "\n" + super.toString();
	}	
}
 		\end{lstlisting}
 	\end{exampleblock}
\end{frame}

\begin{frame}[fragile]{Konkretes Fliegengewicht}
	\begin{exampleblock}{RiverTile}
 		\begin{lstlisting}
public class RiverTile extends TileFlyweight {
	private int foundWater;
	
	public RiverTile(Texture texture) {
		super(texture);
		foundWater = 3;
	}
	
	public void walk(Player player) {
		player.refillWater(foundWater);
	}
	
	@Override
	public String render() {
		return "FLUSS\n" + getTexture().display();
	}
}
 		\end{lstlisting}
 	\end{exampleblock}
\end{frame}

\begin{frame}[fragile]{Konkretes Fliegengewicht}
	\begin{exampleblock}{UnsharedTile}
 		\begin{lstlisting}
public class UnsharedTile extends TileFlyweight {
    private final TileFlyweight tileFlyweight;
    private final int x;
    private final int y;

    public String render() {
        String output = tileFlyweight.render();
        output = output + "Position: (" + x + " / " + y + ")\n";
        return output;
    }

	public void walk(Player player) {
		tileFlyweight.walk(player);
	}
}
 		\end{lstlisting}
 	\end{exampleblock}
\end{frame}

\begin{frame}[fragile]{Fliegengewichtfabrik}
	\begin{exampleblock}{LandscapeTileFactory}
 		\begin{lstlisting}
public class TileFlyweightFactory {
	private Map<LandscapeType, TileFlyweight> landscapeTiles;
	private final Map<String, Texture> textures;
	
	public TileFlyweight createLandscapeTile(LandscapeType key) {
		TileFlyweight landscapeTile = landscapeTiles.get(key);
		
		if(landscapeTile == null) {
			switch (key) {
			case FOREST:
				landscapeTile = new ForestTile(getTexture("Wald"));
				break;
			case DESERT:
				// ...
			}
			landscapeTiles.put(key, landscapeTile);
		}
		return landscapeTile;
	}
        \end{lstlisting}
 	\end{exampleblock}
\end{frame}

\begin{frame}{Allgemeines Klassendiagramm mit Delegation}
    \begin{tikzpicture}[scale=0.7, transform shape]
            \umlclass[x=-5, y=0]{Flyweight}{
            +operation(extrinsicState)
            }{}
            
            \umlclass[x=0, y=3]{ConcreteFlyweight}{
            -intrinsicState
            }{
            +operation(extrinsicState)
            }
            
            \umlclass[x=0, y=-3]{UnsharedConcreteFlyweight}{
            -allState
            }{
            +operation(extrinsicState)
            }
        
            \umlclass[x=5, y=0]{FlyweightFactory}{
            -flyweights : Flyweight[]
            }{
            +getFlyweight(key) : Flyweight
            }
    
            \umlclass[x=5, y=3]{Client}{
            }{
            }
        
            \umlinherit[geometry=-|]{ConcreteFlyweight}{Flyweight}
            \umluniassoc[geometry=-|]{UnsharedConcreteFlyweight}{Flyweight}
            \umlaggreg{FlyweightFactory}{Flyweight}
            \umluniassoc{Client}{FlyweightFactory}
    \end{tikzpicture}
\end{frame}

\frameWithImageOnly{\BILDERPFAD/fliegengewicht2.jpg}

\begin{frame}{Klassendiagramm Wald}
    \begin{tikzpicture}[scale=0.7, transform shape]
        \umlclass[x=-5, y=0]{TreeFlyweight}{
        +operation(extrinsicState)
        }{}
        
        \umlclass[x=0, y=3]{ConcreteTreeFlyweight}{
        -intrinsicState
        }{
        +operation(extrinsicState)
        }
        
        \umlclass[x=0, y=-3]{UnsharedTree}{
        -allState
        }{
        +operation(extrinsicState)
        }
    
        \umlclass[x=5, y=0]{TreeFlyweightFactory}{
        -flyweights : TreeFlyweight[]
        }{
        +getFlyweight(key) : TreeFlyweight
        }

        \umlclass[x=5, y=3]{Forest}{
        }{
        }
    
        \umlimpl[geometry=-|]{ConcreteTreeFlyweight}{TreeFlyweight}
        \umluniassoc[geometry=-|]{UnsharedTree}{TreeFlyweight}
        \umlaggreg{TreeFlyweightFactory}{TreeFlyweight}
        \umluniassoc{Forest}{TreeFlyweightFactory}
    \end{tikzpicture}
\end{frame}

\begin{frame}[fragile]{Fliegengewicht}
	\begin{exampleblock}{TreeFlyweight.java}
 		\begin{lstlisting}
public interface TreeFlyweight {
    void draw(int x, int y);
}
 		\end{lstlisting}
 	\end{exampleblock}
\end{frame}

\begin{frame}[fragile]{Konkretes Fliegengewicht}
	\begin{exampleblock}{ConcreteTreeFlyweight.java}
 		\begin{lstlisting}
public class ConcreteTreeFlyweight implements TreeFlyweight {
	private final String texture;

	public ConcreteTreeFlyweight(String texture) {
		this.texture = texture;
	}

	@Override
	public void draw(int x, int y) {
		System.out.printf("...");
	}
}
 		\end{lstlisting}
 	\end{exampleblock}
\end{frame}

\begin{frame}[fragile]{Nicht geteilter Zustand}
	\begin{exampleblock}{UnsharedTree.java}
 		\begin{lstlisting}
public class UnsharedTree {
	private final TreeFlyweight flyweight;
	private final int x;
	private final int y;
	private final int height;

	public UnsharedTree(TreeFlyweight flyweight, 
            int x, int y, int height) {
		this.flyweight = flyweight;
		this.x = x;
		this.y = y;
		this.height = height;
	}

	public void draw() {
		flyweight.draw(x, y);
		System.out.printf("Hoehe: %d\n", height);
	}
}
 		\end{lstlisting}
 	\end{exampleblock}
\end{frame}

\begin{frame}[fragile]{Fliegengewichtfabrik}
	\begin{exampleblock}{TreeFlyweightFactory}
 		\begin{lstlisting}
public class TreeFlyweightFactory {
	private final Map<String, TreeFlyweight> flyweights;

	public TreeFlyweight getTreeFlyweight(String texture) {
		TreeFlyweight flyweight = flyweights.get(texture);
		if (flyweight == null) {
			flyweight = new ConcreteTreeFlyweight(texture);
			flyweights.put(texture, flyweight);
		}
		return flyweight;
	}
} 		    
 		\end{lstlisting}
 	\end{exampleblock}
\end{frame}

\begin{frame}[fragile]{Client}
	\begin{exampleblock}{Forest}
 		\begin{lstlisting}
public class Forest {
	private TreeFlyweightFactory treeFlyweightFactory;
	private List<UnsharedTree> trees;
	
	public Forest() {
		for (int i = 0; i < 10; i++) {
			if (i % 2 == 0) {
				TreeFlyweight oakTexture =
                    treeFlyweightFactory.getTreeFlyweight("Eiche");
				UnsharedTree oakTree = new UnsharedTree(oakTexture, i, i + 1, i + 30);
				trees.add(oakTree);
			} else if (i % 3 == 0) {
                // ...
			} else {
                // ...
			}
		}
	}
 		\end{lstlisting}
 	\end{exampleblock}
\end{frame}

\begin{frame}[fragile]{Client}
	\begin{exampleblock}{Forest}
 		\begin{lstlisting}
public void drawForest() {
    for (UnsharedTree tree: trees) {
        tree.draw();
    }
}
 		\end{lstlisting}
 	\end{exampleblock}
\end{frame}

\begin{frame}{Delegation Vorteile}
	\begin{itemize}
 		\item Klarheit
 		\item Flexibilität
 		\item Bessere Kontrolle
        \item Keine Vererbung unnötiger Methoden
 	\end{itemize}
\end{frame}
\note{
    Die Delegation hat einige Vorteile gegenüber der Vererbungslösung. Der gemeinsame und der spezifische Zustand sind klar getrennt, was die Wartung und das Verständnis erleichtert. Man kann leicht verschiedene Kombinationen von Flyweight-Objekten und spezifischen Zuständen erstellen, ohne die Klassenhierarchie ändern zu müssen. Da der spezifische Zustand außerhalb des Flyweight-Objekts verwaltet wird, kann er flexibler geändert werden. Und zudem hat man bei Vererbung immer das Problem, dass eventuell Methoden geerbt werden, die man auf diese Art vielleicht nicht haben möchte.
}

\begin{frame}{Delegation Nachteile}
	\begin{itemize}
 		\item Mehr Code
 		\item Zusätzliche Klassen
 	\end{itemize}
\end{frame}
\note{
    Es muss mehr Boilerplate-Code geschrieben werden, um die Delegation zu verwalten und der Einsatz zusätzlicher Klassen wird notwendig. Die erhöhen natürlich die Gesamtkomplexität des Codes.
}

\begin{frame}{Vererbung Vorteile}
	\begin{itemize}
 		\item Einfachheit
 		\item Direkter Zugriff
 		\item Weniger Overhead
 	\end{itemize}
\end{frame}
\note{
    Weniger Klassen sind notwendig, da die nicht-geteilte Klasse selbst das Flyweight-Interface implementiert. Die Klasse hat direkten Zugriff auf sowohl den spezifischen als auch den gemeinsamen Zustand, was die Implementierung vereinfachen kann und es ist nicht notwendig, zusätzliche Delegationsschichten zu schreiben.
}

\begin{frame}{Vererbung Nachteile}
	\begin{itemize}
 		\item Starke Kopplung
 		\item Weniger Modularität
        \item Komplexere Hierarchie
 	\end{itemize}
\end{frame}
\note{
    Der spezifische Zustand wird eng mit dem gemeinsamen Zustand gekoppelt, was die Flexibilität reduzieren kann. Vererbung bedeutet immer starke Kopplung. Änderungen am spezifischen Zustand können eine Änderung in der Flyweight-Klasse erfordern. Durch die Vermischung von gemeinsamem und spezifischem Zustand kann die Codebasis schwieriger zu verstehen und zu warten sein.
}

\begin{frame}{Vergleich: Delegation vs. Vererbung}
    \begin{table}[h]
        \centering
        \begin{tabular}{|l|l|l|}
            \hline
            \textbf{Merkmal} & \textbf{Delegation} & \textbf{Vererbung} \\ \hline
            \textbf{Kopplung} & Niedrig & Höher \\ \hline
            \textbf{Flexibilität} & Hoch & Niedriger \\ \hline
            \textbf{Code-Komplexität} & Höher & Niedriger \\ \hline
            \textbf{Modularität} & Besser & Weniger \\ \hline
            \textbf{Klarheit der Zustände} & Klar getrennt & Vermischt \\ \hline
        \end{tabular}
    \end{table}
    \pause
    \begin{exampleblock}{In den meisten Fällen}
        Delegation bevorzugen!
    \end{exampleblock}
\end{frame}

\note{
    Delegation bietet eine klarere Trennung von gemeinsamem und spezifischem Zustand und ist flexibler. Es eignet sich gut für Szenarien, in denen der spezifische Zustand häufig geändert oder erweitert werden muss, ohne die Fliegengewicht-Objekte zu beeinflussen. Vererbung ist einfacher und direkter und reduziert die Anzahl der benötigten Klassen. Es kann sinnvoll sein, wenn die Zustände weniger variabel sind und die Kopplung zwischen gemeinsamem und spezifischem Zustand akzeptabel ist.
}

\frameWithImageOnly{\BILDERPFAD/fliegengewicht3.jpg}
\note{
    Im letzten Beispiel simulieren wir eine Textverwaltung, d.h. der Benutzer gibt Texte ein und wir verwalten die Texte ähnlich wie Java in einem String-Pool, um Speicher zu sparen. Einzig die Positionen und die Schriftgrößen unterscheiden sich je nach Text.
}

\begin{frame}{Klassendiagramm Textverwaltung}
    \begin{tikzpicture}[scale=0.7, transform shape]
        \umlclass[x=-5, y=0]{StringFlyweight}{
        }{}
        
        \umlclass[x=0, y=3]{ConcreteStringFlyweight}{
        }{
        }
        
        \umlclass[x=0, y=-3]{UnsharedStringInformation}{
        }{
        }
    
        \umlclass[x=5, y=0]{StringFlyweightFactory}{
        }{
        }

        \umlclass[x=5, y=3]{Dialog}{
        }{
        }
    
        \umlimpl[geometry=-|]{ConcreteStringFlyweight}{StringFlyweight}
        \umluniassoc[geometry=-|]{UnsharedStringInformation}{StringFlyweight}
        \umlaggreg{StringFlyweightFactory}{StringFlyweight}
        \umluniassoc{Forest}{StringFlyweightFactory}
    \end{tikzpicture}
\end{frame}

\begin{frame}[fragile]{Fliegengewicht}
	\begin{exampleblock}{StringFlyweight.java}
 		\begin{lstlisting}
public interface StringFlyweight {
    void display();
}
 		\end{lstlisting}
 	\end{exampleblock}
\end{frame}

\begin{frame}[fragile]{Konkretes Fliegengewicht}
	\begin{exampleblock}{ConcreteStringFlyweight.java}
 		\begin{lstlisting}
public class ConcreteStringFlyweight implements StringFlyweight {
    private final String string;

    public ConcreteStringFlyweight(String string) {
        this.string = string;
    }

    @Override
    public void display() {
        System.out.printf("String: %s", string);
    }
}
 		\end{lstlisting}
 	\end{exampleblock}
\end{frame}

\begin{frame}[fragile]{Nicht geteilter Zustand}
	\begin{exampleblock}{UnsharedStringInformation.java}
 		\begin{lstlisting}
public class UnsharedStringInformation {
	private int x;
	private int y;
	private int fontSize;
	private StringFlyweight stringFlyweight;
	
	public UnsharedStringInformation(
                StringFlyweight stringFlyweight, 
			int x, int y, int fontSize) {
		this.stringFlyweight = stringFlyweight;
		this.x = x;
		this.y = y;
		this.fontSize = fontSize;
	}
	
	public void display() {
		stringFlyweight.display();
	}
}
 		\end{lstlisting}
 	\end{exampleblock}
\end{frame}

\begin{frame}[fragile]{Fliegengewichtfabrik}
	\begin{exampleblock}{StringFlyweightFactory}
 		\begin{lstlisting}
public class StringFlyweightFactory {
    private final Map<String, StringFlyweight> flyweights;

    public StringFlyweight getStringFlyweight(String string) {
    	StringFlyweight flyweight = flyweights.get(string);
        if (flyweight == null) {
            flyweight = new ConcreteStringFlyweight(string);
            flyweights.put(string, flyweight);
        }
        return flyweight;
    }

    public int getPoolSize() {
        return flyweights.size();
    }
} 		    
 		\end{lstlisting}
 	\end{exampleblock}
\end{frame}
